% TOP_PROBLEM: {"id": 4600011, "title": "Beta functions", "category_id": 11, "category": {"id": 11, "name": "dynamical_systems", "display_name": "Dynamical Systems", "description": "Problems about long-term behavior of deterministic systems, Hamiltonian dynamics, and periodic orbits.", "slug": "dynamical-systems", "order_index": 11, "created_at": "2026-07-31T15:26:25.671Z"}, "status": "partially_solved", "problem_number": "AMR-045-0011", "set_id": 13, "statusQualification": "Both the usual Perron-branch convention and the source’s reciprocal irreducible-polynomial convention are explicitly specified."}
% FIRST_POSTED: 2026-10-02
% ENTRY_KIND: statement_only
% UnsolvedMath ID 4600011; source catalogue code AMR-045-0011
% !TeX program = lualatex
% Definition and sources only; this file is not an LLM solution attempt.
\documentclass[11pt,a4paper]{article}
\usepackage[margin=25mm]{geometry}
\usepackage{fontspec}
\setmainfont{lmroman10-regular.otf}[BoldFont=lmroman10-bold.otf,
 ItalicFont=lmroman10-italic.otf,BoldItalicFont=lmroman10-bolditalic.otf]
\usepackage{amsmath,amssymb,mathtools}
\usepackage{unicode-math}
\setmathfont{latinmodern-math.otf}
\AtBeginDocument{\let\setminus\smallsetminus}
\usepackage{xurl,hyperref}
\hypersetup{colorlinks=true,urlcolor=blue}
\setcounter{secnumdepth}{0}
\setlength{\parindent}{0pt}
\setlength{\parskip}{5pt}
\setlength{\emergencystretch}{3em}
\begin{document}
\section*{Beta functions}\label{4600011}
{\small UnsolvedMath ID 4600011 · AMR-045-0011 · Dynamical Systems · AMR Open Problem Lists}
\subsection{Definitions and mathematical statement}
Let $P$ be a finite primitive stochastic matrix. For real $t$, let
$P^{[t]}$ have entry $p_{ij}^t$ wherever $p_{ij}>0$ and zero otherwise.
Its Perron eigenvalue
\[
 \beta_P(t)=\rho(P^{[t]})>0
\]
is a real-analytic function of $t$, with $\beta_P(1)=1$.
This distinguished eigenvalue is the usual Perron-branch form of the
beta invariant of the mixing Markov shift presented by $P$.

The source packages the invariant algebraically. Choose multiplicatively
independent positive generators $g_1,\ldots,g_s$ for the group generated
by the nonzero entries of $P$. Express each entry as a Laurent monomial
in $u_j=g_j^t$ and factor
$q(z)=\det(zI-P^{[t]})$ over the fraction field of the resulting Laurent
polynomial ring. Let $b(z)$ be the monic irreducible factor whose root
branch is $\beta_P(t)$. The source calls $1/b(z)$ the beta function;
this specifies its reciprocal-polynomial convention.

Characterize the beta invariants arising from all finite primitive
stochastic matrices. Equivalently, give intrinsic conditions on the
distinguished Perron branch with its algebraic relation, together with
a realization by such a matrix. Matrix size and allowed positive
probabilities are unrestricted. The full characteristic polynomial is
not part of this invariant: irrelevant factors are discarded. A
description limited to entropy, to $\beta_P(0)$, or to matrices with
rational entries would answer a more restricted question.
\subsection{Short English statement}
Characterize the beta invariant carried by the distinguished Perron eigenvalue of a powered stochastic matrix.
\subsection{Sources}
\begin{itemize}
\item[S1] ULAM.AI and the UnsolvedMath Contributors, supplied problems.json, record 4600011 (AMR-045-0011), AMR Open Problem Lists. Catalogue identity and source transcription; CC BY 4.0.
\url{https://huggingface.co/datasets/ulamai/UnsolvedMath}.
\item[S2] Boyle - Open Problems in Symbolic Dynamics (2008), Problem/Question/Conjecture 10.2. Original source indexed by AMR-045-0011.
\url{https://www.math.umd.edu/~mboyle/papers/openfinalsub3nov2007.pdf}.
\end{itemize}
\end{document}
